% GENERATED FILE. Edit canonical lessons, then run npm run generate:books.
% canonical-source-hash: fnv1a64:2a76ccc24e58d99a
% canonical-lessons: MW-C103-nhavno, MW-C103-pakavno, MW-C103-banavno, MW-C103-torno, MW-C103-jorno

\chapter{To bathe, To cook, To make, To break, To join}
\label{ch:mw-to-bathe-to-cook-to-make-to-break-to-join}
\hlchaptermodality{103}

\begin{chapteropening}
\textbf{By the end of this chapter:} \emph{I can say to bathe, to cook, to make, to break and to join.}

Complete the last lesson of chapter 103: i can say to bathe, to cook, to make, to break and to join.
\end{chapteropening}

\section[\texorpdfstring{\mw{न्हावणो} (nhāvṇo)}{न्हावणो (nhāvṇo)}]{\mw{न्हावणो} (nhāvṇo) — to bathe}

\label{lesson:MW-C103-nhavno}



\begin{quote}
Before the new one: say the Marwadi for someone, then the Marwadi for each.
\end{quote}

\subsection*{You'll want to know: \mw{न्हावणो}}
\textbf{\mw{न्हावणो}} — \emph{nhāvṇo} — \textquotedblleft{}to bathe\textquotedblright{}.

\begin{grammarlens}[title={What you've built}]
One more word for this chapter.
\end{grammarlens}

\subsection*{Your turn}
\begin{itemize}
\raggedright
  \item \emph{Say it:} \emph{nhāvṇo}
  \item \emph{Say it:} \emph{nhāvṇo}, once more
  \item \emph{Say it:} say \emph{nhāvṇo}
  \item \emph{Recall:} say the Marwadi for someone, then the Marwadi for each, then say \emph{nhāvṇo} again
\end{itemize}

\begin{tcolorbox}[breakable,colback=teal!4,colframe=teal!35!black,title={Before you move on}]
What does \mw{न्हावणो} mean? (\textquotedblleft{}To bathe\textquotedblright{}.) Say it once more.
\end{tcolorbox}
\section[\texorpdfstring{\mw{पकावणो} (pakāvṇo)}{पकावणो (pakāvṇo)}]{\mw{पकावणो} (pakāvṇo) — to cook}

\label{lesson:MW-C103-pakavno}



\begin{quote}
Before the new one: say the Marwadi for each, then the Marwadi for to bathe.
\end{quote}

\subsection*{You'll want to know: \mw{पकावणो}}
\textbf{\mw{पकावणो}} — \emph{pakāvṇo} — \textquotedblleft{}to cook\textquotedblright{}.

\begin{grammarlens}[title={What you've built}]
One more word for this chapter.
\end{grammarlens}

\subsection*{Your turn}
\begin{itemize}
\raggedright
  \item \emph{Say it:} \emph{pakāvṇo}
  \item \emph{Say it:} \emph{pakāvṇo}, once more
  \item \emph{Say it:} say \emph{pakāvṇo}
  \item \emph{Recall:} say the Marwadi for each, then the Marwadi for to bathe, then say \emph{pakāvṇo} again
\end{itemize}

\begin{tcolorbox}[breakable,colback=teal!4,colframe=teal!35!black,title={Before you move on}]
What does \mw{पकावणो} mean? (\textquotedblleft{}To cook\textquotedblright{}.) Say it once more.
\end{tcolorbox}
\section[\texorpdfstring{\mw{बणावणो} (baṇāvṇo)}{बणावणो (baṇāvṇo)}]{\mw{बणावणो} (baṇāvṇo) — to make}

\label{lesson:MW-C103-banavno}



\begin{quote}
Before the new one: say the Marwadi for to bathe, then the Marwadi for to cook.
\end{quote}

\subsection*{You'll want to know: \mw{बणावणो}}
\textbf{\mw{बणावणो}} — \emph{baṇāvṇo} — \textquotedblleft{}to make\textquotedblright{}.

\begin{grammarlens}[title={What you've built}]
One more word for this chapter.
\end{grammarlens}

\subsection*{Your turn}
\begin{itemize}
\raggedright
  \item \emph{Say it:} \emph{baṇāvṇo}
  \item \emph{Say it:} \emph{baṇāvṇo}, once more
  \item \emph{Say it:} say \emph{baṇāvṇo}
  \item \emph{Recall:} say the Marwadi for to bathe, then the Marwadi for to cook, then say \emph{baṇāvṇo} again
\end{itemize}

\begin{tcolorbox}[breakable,colback=teal!4,colframe=teal!35!black,title={Before you move on}]
What does \mw{बणावणो} mean? (\textquotedblleft{}To make\textquotedblright{}.) Say it once more.
\end{tcolorbox}
\section[\texorpdfstring{\mw{तोड़णो} (toṛṇo)}{तोड़णो (toṛṇo)}]{\mw{तोड़णो} (toṛṇo) — to break}

\label{lesson:MW-C103-torno}



\begin{quote}
Before the new one: say the Marwadi for to cook, then the Marwadi for to make.
\end{quote}

\subsection*{You'll want to know: \mw{तोड़णो}}
\textbf{\mw{तोड़णो}} — \emph{toṛṇo} — \textquotedblleft{}to break\textquotedblright{}.

\begin{grammarlens}[title={What you've built}]
One more word for this chapter.
\end{grammarlens}

\subsection*{Your turn}
\begin{itemize}
\raggedright
  \item \emph{Say it:} \emph{toṛṇo}
  \item \emph{Say it:} \emph{toṛṇo}, once more
  \item \emph{Say it:} say \emph{toṛṇo}
  \item \emph{Recall:} say the Marwadi for to cook, then the Marwadi for to make, then say \emph{toṛṇo} again
\end{itemize}

\begin{tcolorbox}[breakable,colback=teal!4,colframe=teal!35!black,title={Before you move on}]
What does \mw{तोड़णो} mean? (\textquotedblleft{}To break\textquotedblright{}.) Say it once more.
\end{tcolorbox}
\section[\texorpdfstring{\mw{जोड़णो} (joṛṇo)}{जोड़णो (joṛṇo)}]{\mw{जोड़णो} (joṛṇo) — to join}

\label{lesson:MW-C103-jorno}



\begin{quote}
Before the new one: say the Marwadi for to make, then the Marwadi for to break.
\end{quote}

\subsection*{You'll want to know: \mw{जोड़णो}}
\textbf{\mw{जोड़णो}} — \emph{joṛṇo} — \textquotedblleft{}to join\textquotedblright{}.

\begin{grammarlens}[title={What you've built}]
One more word for this chapter.
\end{grammarlens}

\subsection*{Your turn}
\begin{itemize}
\raggedright
  \item \emph{Say it:} \emph{joṛṇo}
  \item \emph{Say it:} \emph{joṛṇo}, once more
  \item \emph{Say it:} say \emph{joṛṇo}
  \item \emph{Recall:} say the Marwadi for to make, then the Marwadi for to break, then say \emph{joṛṇo} again
\end{itemize}

\begin{tcolorbox}[breakable,colback=teal!4,colframe=teal!35!black,title={Before you move on}]
What does \mw{जोड़णो} mean? (\textquotedblleft{}To join\textquotedblright{}.) Say it once more.
\end{tcolorbox}
